\documentclass[10pt]{article}
\usepackage[margin=0.75in,top=0.7in,bottom=0.7in]{geometry}
\usepackage{amsmath,amssymb}
\usepackage{hyperref}
\usepackage{enumitem}
\usepackage{parskip}
\setlength{\parskip}{3pt plus 1pt minus 1pt}

\hypersetup{colorlinks=true,linkcolor=blue,citecolor=blue,urlcolor=blue}


\date{April 14, 2026}

\begin{document}

\begin{center}
    {\large \textbf{Undetectable Watermarking for Autoregressive Image Models: Does the CGZ Impossibility Transfer?}} \\[0.4em]
    % {\large Precise Definitions, Method, and Pixel vs.~Latent Comparison} \\[0.8em]
    {\small Raymond Bahng (\texttt{reugene@mit.edu}), Isabella Wu (\texttt{isawu888@mit.edu}), Maureen Zhang (\texttt{maureenz@mit.edu})} \\[0.4em]
    {\small April 14, 2026}
\end{center}
\thispagestyle{empty}

\paragraph{Background and Central Question:}
Christ, Gunn, and Zamir (CGZ)~[1] construct an undetectable watermarking
scheme for autoregressive language models and prove a fundamental
impossibility (Section~6.2) that for any undetectable scheme on a prefix-specifiable model, a polynomial-time adversary can remove the watermark by re-generating the response one token at a time, since undetectability guarantees each conditional next-token distribution is
statistically close to the clean model's.

Our \textbf{central question} is: \emph{Can we define a weakened notion of undetectability, one that still ensures watermarked outputs are of "good quality" defined by some efficiently computable metric, under which Theorem~11 of CGZ [1] no longer applies?} We study this in the setting of autoregressive (AR) image models such as LlamaGen~[2], which sample discrete visual tokens sequentially. The key distinction from text is
that visual tokens must be decoded through a VQ-VAE to produce pixels. An adversary attempting the CGZ removal attack faces a lossy VQ-VAE round-trip that may introduce visible artifacts, suggesting that successful removal necessarily degrades perceptual quality. This 
motivates a relaxed notion of undetectability in which 
indistinguishability is evaluated under perceptual metrics (FID, LPIPS) rather than total variation distance, and which may be both achievable and sufficient to resist practical removal.

\paragraph{Related work.}
For language models, Christ, Gunn, and Zamir~[1] construct an undetectable watermark and prove in Section 6.2 that any such scheme is inherently removable
by a polynomial-time adversary who re-samples tokens autoregressively. For generative image models, Gunn, Zhao, and Song~[7] present the first provably
undetectable watermark, using a pseudorandom error-correcting code (PRC) construction~[8] to select initial diffusion latents. They demonstrate quality
preservation and robustness on a standard diffusion model. However, their scheme
targets diffusion models, where randomness enters as continuous initial
noise, and does not address AR image models, nor does it engage with the CGZ
impossibility. Recent works~[3,4,5] apply heuristic watermarking to AR models such as 
LlamaGen without formal undetectability guarantees or any analysis of 
the CGZ removal argument. Zhao et al.~[6] study watermark removal via generative AI, but targeting pixel-level perturbations rather than the token-resampling attack we study. Our work is the first to
formally examine CGZ undetectability and the CGZ Section~6.2 impossibility [1] for AR
image models, identifying the VQ-VAE decode-encode path as a structural asymmetry with no text analog that may constrain the attack's effectiveness.

\paragraph{Component 1: Definitional.}
We adapt the CGZ undetectability game to the AR image setting and
propose a formal relaxation suited to this domain. Concretely, we
replace the total-variation distinguishing game with a
\emph{perceptual undetectability} game, where no efficient adversary can
distinguish watermarked from unwatermarked images under any efficiently
computable perceptual metric (e.g.\ an LPIPS-based distinguisher). We then
ask whether this relaxation evades the CGZ removal theorem in Section 6.2 [1]. We
examine two necessary conditions for the attack to go through: (i)
whether LlamaGen supports prefix-specifiability and (ii)~whether
the adversary's regeneration via VQ-VAE decode--re-encode incurs
distortion that the relaxed definition can detect. If the impossibility
applies unchanged, we characterize the 
quality cost of pixel-space regeneration relative to token-space 
re-sampling and formalize any measurable gap as a proof sketch.

\paragraph{Component 2: Empirical.}
We implement CGZ watermarking on LlamaGen's sampling loop for class-conditional ImageNet generation (the scheme is token-semantic-agnostic and ports directly). We measure detection accuracy (TPR/FPR), FID degradation, and per-token entropy across spatial positions, characterizing where the CGZ entropy condition ($\geq\!\sqrt{L}$) holds and where it fails. We then red-team with three attacks:
\begin{enumerate}[leftmargin=*,nosep,topsep=2pt]
  \item \textbf{Token-by-token regeneration}: the Section~6.2 attack [1] applied to LlamaGen, testing feasibility (cost, prefix-specifiability) in the image setting.
  \item \textbf{VQ-VAE decode--re-encode}: pass the watermarked image through the VQ-VAE encoder to obtain perturbed tokens. This image-specific attack has no text analog and tests whether the lossy round-trip removes the watermark and at what quality cost.
  \item \textbf{Diffusion regeneration}: add noise and reconstruct with a separate diffusion model, testing cross-model removal.
\end{enumerate}
For each attack we report watermark survival rate versus quality degradation (LPIPS, FID). If every successful removal strategy degrades quality beyond a measurable threshold, this constitutes empirical evidence that the CGZ removal theorem [1] does not transfer cleanly to the image setting. Together, the two components of our project ask whether a perceptually relaxed notion of undetectability can evade the CGZ removal theorem 
in the AR image setting.

\vspace{0.2em}
\noindent\textbf{References}\par\vspace{0.1em}
\noindent{\footnotesize
[1]~Christ, Gunn, Zamir, ``Undetectable Watermarks for Language Models,'' \textit{COLT}, 2024.\par
[2]~Sun et al., ``Autoregressive Model Beats Diffusion: Llama for Scalable Image Generation,'' 2024.\par
[3]~Jovanovi\'{c} et al., ``Watermarking Autoregressive Image Tokens,'' 2025.\par
[4]~Wen et al., ``Tree-Ring Watermarks,'' \textit{NeurIPS}, 2023.\par
[5]~Fernandez et al., ``Stable Signature,'' \textit{ICCV}, 2023.\par
[6]~Zhao et al., ``Invisible Image Watermarks Are Provably Removable Using Generative AI,'' \textit{NeurIPS}, 2024.\par
[7]~Gunn, Zhao, Song, ``An Undetectable Watermark for Generative Image Models,'' \textit{ICLR}, 2025.\par
[8]~Christ, Gunn, ``Pseudorandom Error-Correcting Codes,'' \textit{CRYPTO}, 2024.\par
}
\end{document}